% !TEX program = pdflatex
\documentclass[12pt,a4paper]{article}

\usepackage[T1]{fontenc}
\usepackage[margin=1in]{geometry}
\usepackage{xcolor}
\usepackage{hyperref}
\usepackage{setspace}
\usepackage{parskip}
\usepackage{enumitem}

\hypersetup{
  colorlinks=true,
  urlcolor=blue!70!black
}

\onehalfspacing
\pagestyle{empty}

\begin{document}

\begin{center}
{\Large\bfseries Personal Statement}\\[4pt]
{\large Olusegun Israel Odeyemi}\\[2pt]
{\normalsize UK Global Talent Visa --- Digital Technology (Exceptional Promise)}
\end{center}

\bigskip

I am a Computer Science and Engineering graduate, Security Engineer, and Technical Founder based in Derby, applying under the Global Talent route in Digital Technology (Exceptional Promise). My work sits where cyber security, applied machine learning, decentralised finance, and regulatory technology meet. Its central output is AMTTP: a compliance platform combining AI risk scoring, smart contracts, privacy-preserving verification, and enterprise compliance workflows.

Building on a Bachelor of Technology in Computer Science and Engineering, my earlier professional experience in cybersecurity, enterprise API integrations, logistics operations, and vendor risk exposed the practical gap that became the focus of AMTTP: financial and digital-asset systems were becoming faster and more distributed, while compliance controls remained fragmented, reactive, and difficult to verify.

To address this, I spent my first three years in the United Kingdom focused on deep, independent research and development. Operating outside traditional institutional structures allowed me to approach the problem from first principles. This period served as the incubation phase for AMTTP and Blind-Spot Decomposition Theory (BSDT). I architected the mathematical framework, machine-learning risk engine, smart-contract enforcement layer, and security model, actively underwriting the platform's development as an independent technical founder.

This focused R\&D phase directly enabled the succession of technical milestones achieved from late 2025 onwards. The recent deployment of the platform, the GitHub codebase, ML benchmarks, security audit artefacts, research preprints, independent academic validation, and a UK patent filing represent the public release of work that systematically matured during those preceding years. My application relies entirely on this verifiable evidence: what I have mathematically proven, built, tested, deployed, published, and had independently recognised.

\textbf{Why the United Kingdom, and why Derby}

The United Kingdom is the right place to develop this work because of its fintech, cyber-security, and regulatory technology ecosystems, and because the FCA's digital-asset agenda makes compliance infrastructure commercially relevant. Derby gives me access to the Midlands' growing technology corridor while keeping me close enough to London for sector engagement. I intend to remain in the UK long term.

\textbf{What I have built: AMTTP}

AMTTP, the Anti-Money Laundering Transaction Trust Protocol, is a blockchain-native compliance system designed to make risk decisions before transaction finality. Most digital-asset compliance tooling operates after settlement, when funds have already moved. AMTTP addresses that gap by combining machine-learning risk scoring, graph-based relationship analysis, sanctions screening, and policy evaluation inside the Ethereum confirmation window, so a transaction can be approved, escrowed, reviewed, or blocked before settlement.

The platform is a working, publicly accessible Sepolia-backed environment that I designed and built end to end. It includes a Flutter consumer application, a Next.js enterprise War Room dashboard, a FastAPI compliance orchestrator connected to nine specialist microservices, a TypeScript Express oracle gateway, 38 AMTTP-authored Solidity files, and TypeScript and Python SDKs. The work has been carried out independently, without institutional funding, formal academic supervision, or venture backing.

\textbf{Applied machine learning and security}

I trained and evaluated knowledge-distilled student models on 625{,}168 Ethereum address samples, including 372 confirmed fraud cases. The production CPU-optimised XGBoost student records ROC-AUC 0.9898; the LightGBM student reaches ROC-AUC 0.9999997 with F1 0.9906; and the ensemble meta-learner reports ROC-AUC 0.9999998 with PR-AUC 0.9997. Distillation lets the system run on commodity CPUs rather than specialist GPU infrastructure, supporting commercial deployment.

I also designed AMTTP to be privacy-preserving and secure-by-design. Through zkNAF, the system supports sanctions non-membership verification, risk proofs, and KYC-related confirmations without unnecessary data exposure, using Groth16 SNARK verification. The contract suite is supported by Slither static analysis, Echidna property-based fuzzing at 50{,}000 iterations with 100-step sequences, and Hardhat integration tests covering deployment, transfer logic, policy enforcement, and cross-contract behaviour.

\textbf{Original research}

Alongside the platform, I produced four preprints: the AMTTP architecture paper on TechRxiv, BSDT, Universal Detection Principle, and The Geometry of System Collapse on Zenodo. These publications are recent, so I do not present them as long-established citation impact. Their value is evidenced through DOI records, public access metrics, supporting code or benchmark evidence, and early independent validation. The first two papers have together attracted 365 views and 210 downloads as of 5 May 2026, without institutional promotion.

BSDT developed from the same research period that produced AMTTP. It diagnoses classifier blind spots without retraining the underlying model, with recall improvements of up to 84.4 percentage points and a 99.94\% false-positive reduction on the Shuttle benchmark. In April 2026, three of my papers were adopted as recommended readings in the postgraduate course CPE 604 by Dr Sunday Adeola Ajagbe, an independent academic with no prior supervisory relationship to me. This is early academic validation rather than mature citation impact, and I present it on that basis.

\textbf{Intellectual property and commercial intent}

On 4 March 2026 I filed UK patent application No.\ 1026066039, an eight-subsystem filing covering machine-learning risk scoring, blockchain compliance architecture, zero-knowledge identity workflows, cross-chain propagation, graph anomaly detection, adaptive friction, policy automation, and system-stability analysis. AMTTP is positioned as pre-settlement Compliance-as-a-Service infrastructure for regulated digital-finance institutions.

\textbf{Why this profile fits Exceptional Promise}

The relevant question for this route is not whether I have followed a conventional UK developer employment path, but whether I show credible potential to become a leading contributor in UK digital technology. I believe the evidence supports that conclusion. I have a Computer Science and Engineering foundation, prior experience in security and enterprise systems, and a public body of independent technical work that includes a deployed platform, a 361-commit full-stack repository, verified smart contracts, machine-learning benchmarks, security audit artefacts, preprints, academic validation, and a UK patent filing.

My earlier commercial and security experience helped me identify the problem; my UK period allowed the research and architecture to mature; and the public evidence from late 2025 onwards shows the point at which the work became ready for deployment, publication, and external review. I am seeking endorsement so that I can continue this work in the United Kingdom: advancing AMTTP towards regulated deployment, sustaining the associated research programme, and contributing to the UK's leadership in secure, intelligent, and commercially relevant digital infrastructure.

\end{document}

